% ============================================================
% CHUONG 1: CO SO LY THUYET VE THUAT TOAN DI TRUYEN
% (Da chinh sua van phong hoc thuat)
% ============================================================
\chapter{Cơ sở lý thuyết về Thuật toán Di truyền}
\label{chap:co-so-ly-thuyet}

Chương này trình bày các nền tảng lý thuyết cần thiết của Thuật toán Di
truyền. Nội dung chương được tổ chức thành bốn
phần chính: mục đầu tiên phát biểu bài toán tối ưu hóa tổng quát và giới
thiệu họ thuật toán tiến hóa; mục thứ hai trình bày cơ sở sinh học cũng
như nguyên lý hoạt động của GA; mục thứ ba phân tích chi tiết các thành
phần và toán tử cốt lõi; mục cuối cùng tổng hợp quy trình tổng thể của
thuật toán và giới thiệu Lập trình Di truyền như một mở rộng trực tiếp
của GA.

\section{Tổng quan về Bài toán Tối ưu và Họ Thuật toán Tiến hóa}
\label{sec:tong-quan}

\subsection*{Bài toán Tối ưu hóa và Không gian Tìm kiếm}
\label{subsec:bai-toan-toi-uu}

Bài toán tối ưu hóa tổng quát có thể được phát biểu như sau: tìm một
vector biến quyết định $x = (x_1, x_2, \ldots, x_n)$ sao cho hàm mục
tiêu $f(x)$ đạt giá trị cực tiểu hoặc cực đại, trong khi $x$ vẫn thỏa
mãn một tập ràng buộc xác định miền khả thi $\Omega$:
\begin{equation}
\min_{x \in \Omega} f(x) \quad \text{hoặc} \quad \max_{x \in \Omega} f(x).
\label{eq:bai-toan-toi-uu}
\end{equation}

Miền $\Omega \subseteq \mathbb{R}^n$, còn được gọi là không gian tìm
kiếm, thường được xác định bởi một hệ ràng buộc đẳng thức và bất đẳng
thức trên các biến quyết định. Bài toán được phân loại là \textbf{có
ràng buộc} khi $\Omega$ là một tập con thực sự của không gian biến, và
\textbf{không ràng buộc} khi $\Omega$ trùng với toàn bộ không gian biến.

Trong thực tiễn, tối ưu hóa thường đối mặt với hai thách thức chủ yếu
như sau.

\begin{itemize}
    \item \textbf{Vấn đề cực trị cục bộ.} Khi hàm mục tiêu không lồi hoặc
    đa cực trị, các thuật toán chỉ khai thác thông tin cục bộ quanh điểm
    hiện tại, chẳng hạn đạo hàm, có xu hướng hội tụ về cực trị cục bộ gần
    với điểm khởi tạo thay vì cực trị toàn cục trên toàn miền $\Omega$.
    Đây chính là hạn chế cốt lõi của các phương pháp tối ưu dựa trên
    gradient khi áp dụng cho bài toán không lồi.
    \item \textbf{Độ phức tạp tổ hợp của bài toán NP-hard.} Đối với
    nhiều lớp bài toán tối ưu tổ hợp, tiêu biểu như bài toán người bán
    hàng hay bài toán lập lịch, số
    lượng phương án khả thi tăng theo hàm mũ hoặc giai thừa của kích
    thước bài toán. Kết quả là, việc duyệt toàn bộ không gian tìm kiếm
    trở nên bất khả thi về mặt tính toán ngay cả với các bài toán có
    kích thước vừa phải.
\end{itemize}

Hai thách thức nêu trên chính là động lực thúc đẩy sự phát triển của
các phương pháp tìm kiếm không dựa trên đạo hàm, đồng thời có khả năng
khám phá nhiều vùng khác nhau của không gian tìm kiếm. Một trong những
hướng tiếp cận tiêu biểu của nhóm phương pháp này là họ thuật toán tiến
hóa, được giới thiệu trong mục tiếp theo.

\subsection*{Họ Thuật toán Tiến hóa}
\label{subsec:ho-thuat-toan-tien-hoa}

Các phương pháp tối ưu cổ điển dựa trên đạo hàm như Gradient Descent hay
phương pháp Newton vận hành trên một điểm duy nhất: ở mỗi bước lặp, điểm
hiện tại được dịch chuyển theo hướng mà đạo hàm chỉ ra. Cách làm này hội
tụ nhanh và có cơ sở lý thuyết chặt chẽ, nhưng đòi hỏi hàm mục tiêu phải
khả vi và chỉ bảo đảm tìm được cực trị cục bộ nằm gần điểm khởi tạo.

Ở một hướng tiếp cận khác, \textbf{metaheuristic} là các chiến lược tìm
kiếm ở mức tổng quát, được xây dựng để làm việc trong chính những điều
kiện mà phương pháp cổ điển không áp dụng được: hàm mục tiêu không khả
vi, có nhiễu, không lồi, hoặc thậm chí không tồn tại biểu thức giải tích
tường minh. Sự đánh đổi ở đây là rõ ràng. Metaheuristic từ bỏ bảo đảm
về tính tối ưu tuyệt đối để nhận lại phạm vi áp dụng rộng hơn nhiều, và
mục tiêu thực tế của nó là đạt được một nghiệm đủ tốt trong khoảng thời
gian chấp nhận được.

Trong số các metaheuristic, nhóm \textbf{tìm kiếm dựa trên quần thể}
duy trì đồng thời nhiều nghiệm ứng viên thay vì một nghiệm duy nhất, rồi
để cả tập hợp đó cùng tiến hóa qua các thế hệ. Việc khảo sát nhiều vùng
của không gian tìm kiếm cùng lúc làm giảm đáng kể nguy cơ hội tụ sớm về
một cực trị cục bộ. Quan trọng hơn, các nghiệm trong quần thể có thể
trao đổi thông tin với nhau để tạo ra nghiệm mới, điều mà một thuật toán
chỉ theo dõi một điểm không thể thực hiện được. \textbf{Họ Thuật toán
Tiến hóa} là đại diện tiêu biểu nhất của
nhóm này, với quá trình tìm kiếm được dẫn dắt bởi nguyên lý chọn lọc tự
nhiên của Darwin.

Mọi thuật toán thuộc họ EA đều vận hành theo một khung chung gồm năm
bước: khởi tạo ngẫu nhiên một quần thể nghiệm, đánh giá chất lượng của
từng nghiệm, ưu tiên những nghiệm tốt hơn khi chọn cá thể sinh sản, sinh
nghiệm mới bằng các toán tử biến đổi ngẫu nhiên, và lặp lại chu trình
cho tới khi thỏa mãn điều kiện dừng. Bốn nhánh chính của họ này khác
nhau chủ yếu ở cách biểu diễn nghiệm và ở toán tử biến đổi được đặt làm
trọng tâm:

\begin{itemize}
    \item \textbf{Thuật toán Di truyền}: đặt
    lai ghép làm toán tử tìm kiếm chính, kết hợp thông tin di truyền từ
    hai cá thể cha mẹ để tạo ra cá thể con. Đây là đối tượng nghiên cứu
    chính của tiểu luận.
    \item \textbf{Chiến lược Tiến hóa}: do
    Rechenberg và Schwefel phát triển cho bài toán tối ưu tham số thực,
    lấy đột biến Gauss làm toán tử chính và bổ sung cơ chế tự thích nghi
    biên độ đột biến trong quá trình tiến hóa.
    \item \textbf{Lập trình Tiến hóa}:
    do Fogel đề xuất, tiến hóa chủ yếu bằng đột biến trên biểu diễn hành
    vi của cá thể và hầu như không sử dụng lai ghép.
    \item \textbf{Lập trình Di truyền}: do
    Koza phát triển, mở rộng GA để tiến hóa trực tiếp các cấu trúc dạng
    cây biểu thức hoặc chương trình máy tính, thường dùng để tự động
    khám phá công thức toán học hoặc chương trình cho một tác vụ đặt ra
    từ trước.
\end{itemize}

Trong phạm vi tiểu luận này, Thuật toán Di truyền được khảo sát trên cả
ba cách mã hóa nghiệm. Mã hóa số thực được dùng cho bài toán tối ưu
không ràng buộc và bài toán tối ưu có ràng buộc, mã hóa hoán vị cho bài
toán người du lịch, còn mã hóa dạng cây biểu thức theo nguyên lý Lập
trình Di truyền được dùng cho bài toán hồi quy ký hiệu. Khung vận hành
năm bước vừa nêu sẽ được cụ thể hóa dần ở các mục tiếp theo, lần lượt
qua cách mã hóa nghiệm, hàm độ thích nghi, và ba nhóm toán tử chọn lọc,
lai ghép và đột biến; riêng nhánh Lập trình Di truyền được trình bày
tách riêng ở Mục~\ref{sec:gp}.


\section{Cơ sở Khoa học và Nguyên lý Hoạt động của GA}
\label{sec:co-so-khoa-hoc}

\subsection*{Nguồn gốc Lịch sử và Nguyên lý Sinh học}
\label{subsec:nguon-goc-lich-su}

Thuật toán Di truyền do John Holland đề xuất năm 1975 tại Đại học
Michigan trong công trình \textit{Adaptation in Natural and Artificial
Systems} [1], và sau đó được David Goldberg hệ thống hóa cũng như phổ
biến rộng rãi hơn thông qua cuốn sách năm 1989 \textit{Genetic
Algorithms in Search, Optimization, and Machine Learning} [2].

Ý tưởng nền tảng của GA là mô phỏng ba nguyên lý cốt lõi trong học
thuyết tiến hóa của Darwin:

\begin{itemize}
    \item \textbf{Chọn lọc tự nhiên}: những cá thể có khả năng thích nghi tốt hơn với môi
    trường sẽ có xác suất sống sót và sinh sản cao hơn.
    \item \textbf{Di truyền}: các đặc điểm của cá thể cha mẹ được truyền
    lại cho thế hệ con thông qua vật chất di truyền.
    \item \textbf{Biến dị}: các sai khác ngẫu nhiên giữa các cá thể,
    hình thành thông qua lai ghép và đột biến, tạo nên tính đa dạng
    di truyền -- điều kiện tiên quyết để quá trình chọn lọc có tác dụng.
\end{itemize}

Dựa trên ba nguyên lý nêu trên, GA thiết lập một phép tương ứng giữa
các khái niệm sinh học và các thành phần thuật toán. Phép tương ứng
này được tóm tắt trong Bảng~\ref{tab:doi-chieu-thuat-ngu}.

\begin{table}[htbp]
\centering
\caption{Bảng đối chiếu thuật ngữ Sinh học và Thuật toán Di truyền}
\label{tab:doi-chieu-thuat-ngu}
\begin{tabular}{lll}
\toprule
\textbf{Sinh học} & \textbf{Thuật ngữ tiếng Anh}
& \textbf{Thuật toán Di truyền} \\
\midrule
Nhiễm sắc thể & Chromosome & Mã hóa nghiệm \\
Gen           & Gene       & Một biến thành phần trong nghiệm \\
Cá thể        & Individual & Một nghiệm ứng viên \\
Quần thể      & Population & Tập hợp các nghiệm tại một thế hệ \\
Độ thích nghi & Fitness    & Giá trị hàm độ thích nghi \\
Thế hệ        & Generation & Một vòng lặp tiến hóa \\
\bottomrule
\end{tabular}
\end{table}

Nhờ phép tương ứng này, một bài toán tối ưu bất kỳ có thể được diễn
dịch thành một quá trình ``tiến hóa nhân tạo''. Cụ thể, thuật toán xuất
phát từ một quần thể các nghiệm ngẫu nhiên; trong suốt quá trình tiến
hóa, các nghiệm có chất lượng tốt hơn sẽ có nhiều cơ hội được chọn
làm cha mẹ hơn. Qua nhiều thế hệ, quần thể dần dịch chuyển về phía
những vùng của không gian tìm kiếm có giá trị hàm mục tiêu tốt.

\subsection*{Phương pháp Mã hóa Nghiệm}
\label{subsec:ma-hoa}

Việc áp dụng GA cho một bài toán cụ thể bắt đầu bằng bước mã hóa: mỗi
nghiệm cần được biểu diễn dưới dạng một ``nhiễm sắc thể'' mà các
toán tử di truyền có thể thao tác trực tiếp. Bốn cách mã hóa phổ biến
nhất được trình bày như sau.

\begin{itemize}
    \item \textbf{Mã hóa nhị phân}: nghiệm được
    biểu diễn bằng một chuỗi bit gồm $0$ và $1$, chẳng hạn
    \texttt{10110010}. Đây là cách mã hóa nguyên thủy nhất của GA cổ
    điển do Holland đề xuất, phù hợp với các bài toán rời rạc hoặc
    logic. Tuy nhiên, khi áp dụng cho biến số thực, cách mã hóa này
    đòi hỏi một bước rời rạc hóa, làm giảm độ chính xác của nghiệm.
    \item \textbf{Mã hóa số thực}: nghiệm là
    một dãy số thực, chẳng hạn $[3.14,\, -0.05,\, 12.8]$, trong đó mỗi
    phần tử tương ứng trực tiếp với một biến quyết định. Cách mã hóa
    này phù hợp một cách tự nhiên với các bài toán tối ưu liên tục
    nhiều biến và loại trừ được sai số lượng tử hóa của mã hóa nhị
    phân. Đây cũng là cách mã hóa được sử dụng trong phần cài đặt thực
    nghiệm của tiểu luận tại Chương~\ref{chap:ung-dung}.
    \item \textbf{Mã hóa hoán vị}: nghiệm là
    một thứ tự sắp xếp của một tập hợp phần tử, chẳng hạn
    $[3, 1, 4, 2, 5]$. Đây là cách mã hóa tự nhiên cho các bài toán
    định tuyến và lập lịch như TSP hay bài toán phân công công việc,
    nơi nghiệm về bản chất là một hoán vị.
    \item \textbf{Mã hóa dựa trên cấu trúc}:
    nghiệm được biểu diễn dưới dạng cây biểu thức hoặc đồ thị, chủ
    yếu được sử dụng trong Lập trình Di truyền nhằm tiến hóa các công
    thức hoặc chương trình.
\end{itemize}

Cần lưu ý rằng việc lựa chọn cách mã hóa quyết định trực tiếp đến loại
toán tử lai ghép và đột biến có thể áp dụng. Vấn đề này sẽ được phân
tích chi tiết trong Mục~\ref{sec:toan-tu-cot-loi}.


\section{Các Thành phần và Phép toán Cốt lõi}
\label{sec:toan-tu-cot-loi}

\subsection*{Hàm độ thích nghi và Ràng buộc}
\label{subsec:ham-thich-nghi}

Hàm độ thích nghi đóng vai trò đánh giá chất lượng của từng cá thể
trong quần thể, và là cơ sở cho toàn bộ quá trình chọn lọc. Đối với bài
toán không ràng buộc, hàm thích nghi thường được xây dựng trực tiếp từ
hàm mục tiêu; chẳng hạn, đối với bài toán cực tiểu hóa, ta có thể đảo
dấu hàm mục tiêu do GA nguyên thủy được phát biểu dưới dạng cực đại
hóa độ thích nghi.

Đối với bài toán \textbf{có ràng buộc}, trọng tâm của tiểu luận này,
cách tiếp cận kinh điển là đưa ràng buộc vào hàm mục tiêu thông qua một
số hạng phạt, tạo thành hàm thích nghi đã phạt:
\begin{equation}
F(x) = f(x) \pm \sum_{i} \lambda_i \cdot g_i(x),
\label{eq:ham-phat}
\end{equation}
trong đó $g_i(x)$ đo mức độ vi phạm ràng buộc thứ $i$ (nhận giá trị $0$
nếu ràng buộc được thỏa mãn), và $\lambda_i > 0$ là hệ số phạt tương
ứng. Dấu của số hạng phạt cũng như độ lớn của $\lambda_i$ được chọn
tùy theo bài toán là cực tiểu hay cực đại hóa, sao cho một cá thể vi
phạm ràng buộc luôn bị đánh giá thấp hơn một cá thể khả thi có cùng
giá trị hàm mục tiêu.

Hàm phạt tĩnh tồn tại một hạn chế cố hữu: hệ số $\lambda_i$ phải được
xác định thủ công và giữ cố định trong suốt quá trình tiến hóa. Nếu
chọn $\lambda_i$ quá nhỏ, quần thể có xu hướng bỏ qua các ràng buộc;
ngược lại, nếu chọn $\lambda_i$ quá lớn, số hạng phạt sẽ lấn át hoàn
toàn hàm mục tiêu, khiến quần thể chỉ tập trung thỏa mãn ràng buộc mà
không còn khả năng phân biệt giữa các nghiệm khả thi tốt và kém. Để
khắc phục hạn chế này, nhiều kỹ thuật xử lý ràng buộc nâng cao đã được
đề xuất, tiêu biểu như quy tắc khả thi của Deb hoặc ngưỡng
$\varepsilon$ giảm dần theo thời gian do Takahama và Sato đề xuất. Các
kỹ thuật này sẽ được trình bày cụ thể cùng với phần cài đặt thực nghiệm
trong Chương 2.

\subsection*{Các Phép chọn lọc}
\label{subsec:chon-loc}

Toán tử chọn lọc quyết định những cá thể nào của quần thể hiện tại sẽ
được chọn làm cha mẹ để sinh ra thế hệ tiếp theo. Nguyên tắc chung của
toán tử này là thiên vị có hướng về phía các cá thể có độ thích nghi
cao. Ba phương pháp chọn lọc phổ biến nhất được trình bày dưới đây.

\begin{itemize}
    \item \textbf{Vòng quay Roulette}: xác
    suất chọn cá thể $i$ tỉ lệ thuận với độ thích nghi của nó so với tổng
    độ thích nghi của toàn quần thể:
    \begin{equation}
    P_i = \frac{f_i}{\sum_{j=1}^{N} f_j}.
    \label{eq:roulette}
    \end{equation}
    Hạn chế của phương pháp này là khi một cá thể có độ thích nghi vượt
    trội hẳn so với phần còn lại, cá thể đó sẽ chiếm phần lớn xác suất
    được chọn, từ đó làm quần thể nhanh chóng mất đi tính đa dạng.
    \item \textbf{Thách đấu}: chọn ngẫu nhiên $k$
    cá thể từ quần thể, sau đó chọn cá thể tốt nhất trong nhóm làm cha
    mẹ. Tham số $k$, gọi là kích thước thách đấu, kiểm soát trực tiếp
    áp lực chọn lọc: $k$ càng lớn thì áp lực chọn lọc càng cao, quần
    thể hội tụ càng nhanh, nhưng đổi lại rủi ro mất đa dạng cũng tăng
    theo.
    \item \textbf{Chọn lọc theo thứ hạng}: cá thể
    được chọn dựa trên thứ hạng của nó trong quần thể đã được sắp xếp
    theo độ thích nghi, thay vì dựa trên giá trị thích nghi tuyệt đối.
    Cách tiếp cận này khắc phục được hiện tượng một cá thể vượt trội áp
    đảo toàn bộ xác suất chọn như trong phương pháp vòng quay Roulette.
\end{itemize}

Bên cạnh các toán tử chọn lọc cha mẹ nêu trên, chiến lược lựa chọn cá
thể chuyển sang thế hệ sau còn bao gồm cơ chế giữ lại cá thể ưu tú,
được trình bày riêng trong Mục~\ref{subsec:elitism} do vai trò và cách
áp dụng của nó khác biệt so với chọn lọc cha mẹ thông thường.

\subsection*{Các Phép lai ghép}
\label{subsec:lai-ghep}

Toán tử lai ghép kết hợp vật chất di truyền của hai cá thể cha mẹ để
tạo ra cá thể con mới. Tần suất lai ghép $p_c$ thường được đặt trong
khoảng từ $0.6$ đến $0.95$, nghĩa là phần lớn các cặp cha mẹ được chọn
sẽ thực sự tham gia lai ghép, trong khi phần còn lại được sao chép
nguyên vẹn sang thế hệ con.

Đối với mã hóa nhị phân và số thực, ba toán tử lai ghép phổ biến nhất
là:
\begin{itemize}
    \item \textbf{Lai ghép đơn điểm}: chọn
    ngẫu nhiên một vị trí cắt trên chuỗi nhiễm sắc thể. Hai cá thể con
    được tạo ra bằng cách hoán đổi các đoạn nằm sau vị trí cắt giữa hai
    cha mẹ.
    \item \textbf{Lai ghép đa điểm}: là dạng
    tổng quát của lai ghép đơn điểm với nhiều vị trí cắt; các đoạn xen
    kẽ giữa hai cha mẹ được hoán đổi cho nhau.
    \item \textbf{Lai ghép đồng nhất}: giá trị của
    mỗi gen tại vị trí con được quyết định độc lập và ngẫu nhiên, lấy
    từ cha hoặc mẹ, không phụ thuộc vào vị trí của gen trong chuỗi.
\end{itemize}

Đối với mã hóa số thực, một họ toán tử quan trọng khác là lai ghép
trộn, cũng chính là toán tử được sử
dụng trong phần cài đặt của tiểu luận. Xét hai cá thể cha mẹ
$p_1, p_2 \in \mathbb{R}^n$ và hệ số trộn $\alpha$ được lấy mẫu ngẫu
nhiên, có thể nhận giá trị nằm ngoài khoảng $[0, 1]$, hai cá thể con
được tạo ra theo tổ hợp tuyến tính:
\begin{equation}
c_1 = \alpha \cdot p_1 + (1-\alpha) \cdot p_2,
\qquad
c_2 = \alpha \cdot p_2 + (1-\alpha) \cdot p_1.
\label{eq:blend-crossover}
\end{equation}
Khi $\alpha$ được lấy mẫu từ một khoảng mở rộng ra ngoài $[0, 1]$,
chẳng hạn $\alpha \sim \mathcal{U}(-0.25,\, 1.25)$, cá thể con có thể
nằm ngoài đoạn thẳng nối $p_1$ và $p_2$. Nhờ vậy, toán tử này cân bằng
được giữa khả năng khai thác (khai thác vùng lân cận cha mẹ) và khả
năng khám phá (tiếp cận các vùng nằm ngoài đoạn nối hai cha mẹ), thay
vì chỉ co dần khoảng tìm kiếm qua các thế hệ.

Đối với mã hóa hoán vị trong các bài toán định tuyến và lập lịch, việc
hoán đổi trực tiếp từng đoạn như các toán tử trên có thể tạo ra cá thể
con vi phạm tính hoán vị do có phần tử bị lặp lại hoặc thiếu sót. Vì
vậy, các toán tử lai ghép chuyên biệt cần được sử dụng:
\begin{itemize}
    \item \textbf{PMX}: hoán đổi một đoạn
    con giữa hai cha mẹ theo cách tương tự lai ghép đơn điểm, sau đó
    điều chỉnh các vị trí còn lại dựa trên ánh xạ tương ứng từng phần
    nhằm đảm bảo tính hoán vị của cá thể con.
    \item \textbf{OX}: giữ nguyên một đoạn con từ cha
    mẹ thứ nhất, các phần tử còn lại được điền vào theo đúng thứ tự xuất
    hiện của chúng trong cha mẹ thứ hai.
    \item \textbf{CX}: xác định các ``chu trình'' giữa
    vị trí phần tử của hai cha mẹ; mỗi cá thể con được tạo ra bằng cách
    giữ nguyên các chu trình xen kẽ lấy từ mỗi cha mẹ.
\end{itemize}

\subsection*{Các Phép đột biến}
\label{subsec:dot-bien}

Toán tử đột biến tạo ra biến dị ngẫu nhiên trên một cá thể. Tần suất
đột biến $p_m$ trên mỗi gen thường nhỏ, rơi vào khoảng từ $0.001$ đến
$0.05$ đối với GA nhị phân cổ điển. Vai trò chính của đột biến là duy
trì đa dạng di truyền cho quần thể, đồng thời giúp thuật toán thoát ra
khỏi những vùng mà toán tử lai ghép một mình không thể tiếp cận được.
Một số dạng đột biến tiêu biểu bao gồm:
\begin{itemize}
    \item \textbf{Đột biến đảo bit}: áp dụng cho mã
    hóa nhị phân; mỗi bit được đảo giá trị ($0 \leftrightarrow 1$) độc
    lập với xác suất $p_m$.
    \item \textbf{Đột biến Gauss}: áp dụng cho mã
    hóa số thực; mỗi biến $x_i$ được cộng thêm một nhiễu ngẫu nhiên, phổ
    biến nhất là nhiễu Gauss:
    \begin{equation}
    x_i' = x_i + \mathcal{N}(0,\, \sigma_i^2),
    \label{eq:dot-bien-gauss}
    \end{equation}
    trong đó độ lệch chuẩn $\sigma_i$ thường được đặt tỉ lệ với độ rộng
    miền giá trị của biến $x_i$, nhằm đảm bảo bước đột biến có ý nghĩa
    tương đương giữa các biến có thang đo khác nhau.
    \item \textbf{Đột biến hoán đổi, đảo ngược và chèn}: áp dụng cho mã hóa hoán vị; lần lượt là hoán
    đổi vị trí hai phần tử, đảo ngược thứ tự một đoạn con, và di chuyển
    một phần tử sang vị trí khác trong chuỗi. Cả ba thao tác đều bảo
    toàn tính hoán vị của cá thể.
\end{itemize}

\subsection*{Chiến lược giữ lại Cá thể Ưu tú}
\label{subsec:elitism}

Do các toán tử chọn lọc, lai ghép và đột biến đều mang tính ngẫu nhiên,
không có cơ chế nào đảm bảo cá thể tốt nhất của thế hệ hiện tại sẽ được
duy trì sang thế hệ kế tiếp. Trên lý thuyết, độ thích nghi tốt nhất của
toàn quần thể hoàn toàn có thể giảm đi giữa hai thế hệ liên tiếp. Chiến
lược \textbf{Elitism} ra đời nhằm khắc phục vấn đề này: $k$ cá thể tốt
nhất của thế hệ hiện tại được giữ nguyên vẹn và chuyển trực tiếp sang
thế hệ sau mà không trải qua lai ghép hay đột biến. Nhờ đó, độ thích
nghi tốt nhất từng đạt được sẽ không bao giờ giảm qua các thế hệ.

Việc lựa chọn kích thước elitism $k$ đòi hỏi một sự cân nhắc. Nếu $k$
quá nhỏ, các cá thể tốt vẫn có thể bị mất do tác động của các toán tử
ngẫu nhiên. Ngược lại, nếu $k$ quá lớn, áp lực khám phá của quần thể sẽ
giảm sút, làm suy giảm tính đa dạng và khiến quá trình tìm kiếm hội tụ
sớm về một vùng hẹp của không gian tìm kiếm.


\section{Quy trình Thuật toán và Điều kiện Dừng}
\label{sec:quy-trinh}

\subsection*{Sơ đồ khối và Giả mã}
\label{subsec:so-do-gia-ma}

Phối hợp các thành phần đã trình bày trong
Mục~\ref{sec:toan-tu-cot-loi}, một thế hệ của GA điển hình bao gồm
chu trình lặp lại sau: chọn lọc cha mẹ, lai ghép, đột biến, đánh giá
quần thể con, rồi cập nhật quần thể mới với cơ chế elitism. Chu trình
này được minh họa tổng quát trong Hình~\ref{fig:ga-cycle}.

\begin{figure}[htbp]
\centering
\begin{tikzpicture}[
    node distance=0.75cm and 2.4cm,
    box/.style={draw, rounded corners, fill=blue!8, align=center,
                minimum width=4.4cm, minimum height=0.9cm, font=\small},
    decision/.style={draw, diamond, aspect=2.2, fill=orange!12,
                align=center, text width=2.3cm, inner sep=1pt, font=\small},
    arrow/.style={-{Latex[length=2mm]}}
]

\node[box] (init) {Khởi tạo quần thể $P(0)$};
\node[box, below=of init] (eval0) {Đánh giá độ thích nghi};
\node[box, below=of eval0] (select) {Chọn lọc (Selection)};
\node[box, below=of select] (cross) {Lai ghép (Crossover)};
\node[box, below=of cross] (mut) {Đột biến (Mutation)};
\node[box, below=of mut] (evalc) {Đánh giá quần thể con};
\node[box, below=of evalc] (update) {Cập nhật quần thể (Elitism)};
\node[decision, below=of update] (stop) {Đạt điều kiện dừng?};
\node[box, below=of stop] (ret) {Trả về cá thể tốt nhất};

\draw[arrow] (init) -- (eval0);
\draw[arrow] (eval0) -- (select);
\draw[arrow] (select) -- (cross);
\draw[arrow] (cross) -- (mut);
\draw[arrow] (mut) -- (evalc);
\draw[arrow] (evalc) -- (update);
\draw[arrow] (update) -- (stop);
\draw[arrow] (stop) -- node[right]{Đúng} (ret);

\draw[arrow] (stop.west) -- ++(-2.3cm,0) coordinate (loopcorner) |-
(select.west);
\node[left] at (loopcorner) {\small Sai, $t = t+1$};

\end{tikzpicture}
\caption{Chu trình tiến hóa của Thuật toán Di truyền}
\label{fig:ga-cycle}
\end{figure}

Để mô tả chính xác hơn, Thuật toán~\ref{alg:ga-pseudocode} trình bày
chu trình tiến hóa trên dưới dạng giả mã.

\begin{algorithm}[htbp]
\caption{Giả mã tổng quát của Thuật toán Di truyền}
\label{alg:ga-pseudocode}
\begin{algorithmic}[1]
\State $t \gets 0$
\State Khởi tạo quần thể ban đầu $P(0)$ ngẫu nhiên
\State Đánh giá độ thích nghi của $P(0)$
\While{điều kiện dừng chưa thỏa mãn}
    \State $t \gets t + 1$
    \State Chọn lọc tập cha mẹ $S(t)$ từ $P(t-1)$
    \State Thực hiện lai ghép (crossover) trên $S(t)$, tạo tập con $C(t)$
    \State Thực hiện đột biến (mutation) trên $C(t)$
    \State Đánh giá độ thích nghi của $C(t)$
    \State $P(t) \gets$ cập nhật quần thể mới (áp dụng elitism)
\EndWhile
\State \Return cá thể tốt nhất tìm được
\end{algorithmic}
\end{algorithm}

\subsection*{Tiêu chuẩn Dừng và Đánh giá sự Hội tụ}
\label{subsec:tieu-chuan-dung}

Một thuật toán tiến hóa cần được trang bị một điều kiện dừng tường
minh để kết thúc quá trình lặp. Ba tiêu chuẩn dừng phổ biến, thường
được kết hợp sử dụng với nhau, được trình bày như sau.

\begin{itemize}
    \item \textbf{Số thế hệ tối đa} $G_{max}$: đây là tiêu chuẩn đơn
    giản nhất, đồng thời giúp kiểm soát trực tiếp ngân sách tính toán.
    Tuy nhiên, việc lựa chọn $G_{max}$ cần cân nhắc: nếu $G_{max}$ quá
    nhỏ, thuật toán có thể dừng lại trước khi quần thể hội tụ; nếu
    $G_{max}$ quá lớn, tài nguyên tính toán sẽ bị lãng phí sau khi quần
    thể đã hội tụ.
    \item \textbf{Ngưỡng độ thích nghi mục tiêu}: thuật toán dừng ngay
    khi tìm được một cá thể có độ thích nghi đạt hoặc vượt ngưỡng đặt
    trước. Tiêu chuẩn này phù hợp khi giá trị tối ưu cần đạt đã biết
    trước hoặc có thể ước lượng được.
    \item \textbf{Độ trễ hội tụ}: thuật toán dừng khi độ
    thích nghi của cá thể tốt nhất không cải thiện sau $K$ thế hệ liên
    tiếp, được xem là dấu hiệu cho thấy các toán tử di truyền không còn
    tạo ra tiến bộ đáng kể. Tuy vậy, cần lưu ý rằng tiêu chuẩn này có
    thể khiến thuật toán dừng quá sớm nếu $K$ quá nhỏ, do quần thể hoàn
    toàn có thể trải qua giai đoạn không cải thiện tạm thời trước khi
    tiếp tục tiến triển ở các thế hệ sau.
\end{itemize}


\section{Lập trình Di truyền}
\label{sec:gp}

Lập trình Di truyền do John Koza phát triển
từ cuối thập niên 1980 và được hệ thống hóa trong công trình năm 1992
[8]. GP là một mở rộng trực tiếp của GA, tương ứng với dạng \textit{mã
hóa dựa trên cấu trúc} đã được giới thiệu trong
Mục~\ref{subsec:ma-hoa}. Khác biệt cốt lõi so với GA dùng mã hóa số
thực hay hoán vị là: thay vì tiến hóa một chuỗi giá trị có độ dài cố
định trong khuôn dạng nghiệm đã cho trước, GP tiến hóa trực tiếp
\textbf{cấu trúc} của một biểu thức hoặc chương trình. Nhờ đó, thuật
toán có khả năng tự khám phá đồng thời cả hình dạng lẫn tham số của
nghiệm.

\subsection*{Mã hóa dạng cây biểu thức}

Mỗi cá thể trong GP là một cây biểu thức. Các nút trong của cây là một
toán tử lấy từ \textit{tập hàm} cho trước, chẳng hạn các phép toán số
học $\{+,-,\times,\div\}$ hoặc hàm lượng giác $\{\sin,\cos\}$; các nút
lá là biến đầu vào hoặc hằng số lấy từ \textit{tập đầu cuối}. Khi duyệt
cây theo đúng cấu trúc phân cấp của nó, ta thu được một biểu thức toán
học tường minh, có thể tính giá trị trên bất kỳ bộ dữ liệu đầu vào nào.
Điều này cho phép toàn bộ quần thể các cây với hình dạng và nội dung
khác nhau được đánh giá, so sánh một cách thống nhất như trong GA
thông thường.

\subsection*{Hàm độ thích nghi}

Ứng dụng tiêu biểu nhất của GP là bài toán \textbf{hồi quy ký hiệu}:
tự động tìm một biểu thức toán học mô tả tốt mối quan hệ giữa tập dữ
liệu quan sát $(X, y)$ mà không giả định trước dạng hàm. Bài toán này
khác biệt căn bản so với hồi quy tuyến tính hay hồi quy đa thức, nơi
dạng hàm đã được ấn định sẵn và công việc còn lại chỉ là tối ưu các hệ
số. Trong GP, độ thích nghi của một cây thường được xây dựng từ sai số
dự đoán, tiêu biểu là sai số bình phương trung bình giữa giá trị
cây tính được và giá trị $y$ thực tế. Ngoài ra, một số hạng phạt tỉ lệ
với kích thước cây thường được cộng thêm,
nhằm hạn chế hiện tượng cây phát triển không kiểm soát
mà không mang lại cải thiện tương xứng về độ chính xác.

\subsection*{Toán tử tiến hóa}

Các toán tử lai ghép và đột biến trong GP thao tác trực tiếp trên cấu
trúc cây thay vì trên chuỗi có độ dài cố định:
\begin{itemize}
    \item \textbf{Lai ghép trao đổi cây con}: chọn
    ngẫu nhiên một nút trên mỗi cây cha mẹ, sau đó hoán đổi hai cây con
    bắt rễ tại các nút đó cho nhau để tạo ra hai cây con mới.
    \item \textbf{Đột biến thay thế cây con}: chọn
    ngẫu nhiên một nút trên cây, rồi thay toàn bộ cây con bắt rễ tại nút
    đó bằng một cây con mới được sinh ngẫu nhiên.
\end{itemize}

Các thành phần còn lại của chu trình tiến hóa, bao gồm khởi tạo quần
thể, chọn lọc, elitism và điều kiện dừng, về bản chất tương tự như GA
đã trình bày trong các mục trước; khác biệt duy nhất nằm ở đối tượng
thao tác là \textbf{cây} thay vì \textbf{vector} số hoặc hoán vị.

Tóm lại, chương này đã trình bày đầy đủ các nền tảng lý thuyết của GA
và GP. Trên cơ sở đó, Chương 2 sẽ trình bày việc áp dụng GA và GP cho
bốn nhóm bài toán: tối ưu không ràng buộc, tối ưu có ràng buộc, bài
toán người du lịch và hồi quy ký hiệu. Kết quả thực nghiệm sẽ được đối
chiếu với các mốc so sánh tương ứng của từng nhóm, gồm giá trị cực tiểu
toàn cục đã biết của các hàm benchmark, nghiệm tối ưu công bố của bộ
chuẩn CEC 2006 và của TSPLIB95, cùng với phương pháp hồi quy tuyến
tính. Độc giả quan tâm đến các hướng phát triển và ứng dụng khác của GA
có thể tham khảo thêm tổng quan của Katoch và cộng sự năm 2020 [5], cũng
như một ứng dụng kinh điển kết hợp GA với huấn luyện mạng nơ-ron của
Whitley và cộng sự năm 1990 [6].
